% this file is encoded in utf-8
% v3.0 (Jun. 11, 2019)

\documentclass[12pt, a4paper]{ntuthesis}

% 除非校方修改了論文格式 (margins, header, footer, 浮水印, 中文數字之章別)
% 或者需要增加所用的 LaTeX 套件，
% 或者要改預設中文字型、編碼
% 否則毋須修改本檔內容

%%%%%%%%%%%%%%%%%%%%%%%%%%%%%%
\usepackage{mathtools} % 各式 AMS 數學功能
\usepackage{amssymb} % 各式 AMS 數學符號
\usepackage{mathrsfs} %草寫體數學符號，在數學模式裡用 \mathscr{E} 得草寫 E
\usepackage{bm}
\usepackage{caption}
\usepackage{subcaption}
\usepackage{tabularx}
\usepackage{url}
\usepackage[usenames,dvipsnames]{color}
%\usepackage[square, comma, numbers, sort&compress]{natbib}
\usepackage[
    colorlinks, 
    linkcolor = black, 
    citecolor = red,    % black
    urlcolor  = blue,   % black
    unicode
]{hyperref}

\usepackage{siunitx}
\usepackage{tikz}
\usetikzlibrary{arrows.meta,positioning,calc,chains,scopes}
\usetikzlibrary{shapes.geometric,angles}
\usetikzlibrary{backgrounds,fit}

\usepackage{multirow,multicol,rotating}
% 插圖套件 graphicx
\usepackage{graphicx}
% font, margin, linestretch 是設在 ntuthesis.sty 裡
\usepackage{ntuthesis}


%%%%%%%%% listings / algorithm %%%%%%%%%%
\usepackage{listings} % 程式列表套件
% hyperref跟algorithm衝突，hyperref必須放在algorithm前面
\usepackage{algorithm}
\usepackage{algorithmic}

%
% listing setting
\lstset{breaklines=true,% 過長的程式行可斷行
extendedchars=false,% 中文處理不需要 extendedchars
texcl=true,% 中文註解需要有 TeX 處理過的 comment line, 所以設成 true
comment=[l]\%\%,% 以雙「百分號」做為程式中文註解的起頭標記，配合 MATLAB
basicstyle=\small,% 小號字體, 約 10 pt 大小
commentstyle=\upshape,% 預設是斜體字，會影響註解裏的英文，改用正體
%language=Octave % 會將一些 octave 指令以粗體顯示
}
%%%%%%%%%%%%%%%%%%%%%%%%%%%%%%%%%%%%%%%%%

%%%%%%%%% fancyhdr %%%%%%%%%%
% 增強功能型頁楣 / 頁腳套件
\usepackage{fancyhdr}  % 借用此套件來擺放浮水印 
% (佔用了 central header)
% 不需要浮水印的使用者仍可利用此套件，產生所需的 header, footer
%
% 啟動 fancy header/footer 套件
\pagestyle{fancy}
\fancyhead{}  % reset left, central, right header to empty
\fancyfoot[C]{\thepage} %中間 footer 擺放頁碼
\renewcommand{\headrulewidth}{0pt} % header 的直線; 0pt 則無線

% 如果不需要任何浮水印，則請把下列介於 >>> 與 <<< 之間
% 的文字行關掉 (行首加上百分號)
%% 浮水印 >>> 
% 請用 Adobe Acrobat 加入
% 也會需要加入 doi 的浮水印
%% this file is encoded in utf-8
% v3.0 (Jun. 11, 2019)

% 如果浮水印不是全篇需要，請把下列介於 >>> 與 <<<
% 的「全篇浮水印專用碼」關掉 (行首加百分號)
% 參考自 Keith Reckdahl 寫的 "Using Imported Graphics in LATEX2e" (epslatex.pdf) p.39
% 如果只有特定頁需要浮水印
% 則依該頁屬性使用下列之一的命令 
% 普通頁命令 \thispagestyle{WaterMarkPage}
% plain 頁命令 \thispagestyle{PlainWaterMarkPage}
% empty 頁命令 \thispagestyle{EmptyWaterMarkPage}


% 將重複使用的浮水印章
% 圖檔是 watermark.xxx
% 副檔名可以不加，可以是 latex 系統能處裡的任何格式：pdf, gif, png, jpg, eps, ...
% 某些圖檔格式在某些工作流程可能需要作前置處裡。
% 例如，pdflatex 無法直接處理 eps 檔
%  latex + dvipdfmx 無法直接處理 pdf, gif, png, jpg, 需要用 ebb 小工具程式
%  對圖檔產生 .bb 對應檔。
%
% 寬為 5.1 cm
\newsavebox{\mywatermark}
\sbox{\mywatermark}{\includegraphics[keepaspectratio, scale=0.5]{watermark}}


% 將 central header 擺放浮水印的巨集指令
\newcommand{\PlaceWaterMark}{\fancyhead[C]{\setlength{\unitlength}{1cm}%
\begin{picture}(0,0)%
\put(4.4,-4.2){\usebox{\mywatermark}}% 圖檔擺放的位置座標
\end{picture}}%
}

\fancyhead{}  % reset left, central, right header to empty
%% 如果不需整篇論文都要浮水印
%% 則下面  >>> 與 <<< 之間的程式碼請關閉
%% >>> 全篇浮水印
\PlaceWaterMark  % 每一頁都有浮水印 (除了 plain、empty 頁以外)

% 重新定義 plain 頁面
% 每張 plain 頁面 (每一章的第一頁) 也加浮水印

\fancypagestyle{plain}{%
\fancyhead{}%
\PlaceWaterMark%
\fancyfoot{}%
\fancyfoot[C]{\thepage}
\renewcommand{\headrulewidth}{0pt}%
\renewcommand{\footrulewidth}{0pt}%
}
%% <<< 全篇浮水印

%% 如果只有一、兩頁需要有浮水印
%% 可以在該頁 (有頁碼) 使用 \thispagestyle{WaterMarkPage}
%% 此命令不影響原有的 header、footer
\fancypagestyle{WaterMarkPage}{%
\PlaceWaterMark%
}

%% 如果只有一、兩頁 plain 頁需要有浮水印 (如 摘要、自傳等)
%% 可以在該頁 (有頁碼) 使用 \thispagestyle{PlainWaterMarkPage}
%% 只有頁碼與浮水印，沒有其他的 header、footer
%% 等同於 plain page style + water mark
\fancypagestyle{PlainWaterMarkPage}{%
\fancyhead{}%
\PlaceWaterMark%
\fancyfoot{}%
\fancyfoot[C]{\thepage}
\renewcommand{\headrulewidth}{0pt}%
\renewcommand{\footrulewidth}{0pt}%
}

%% 如果只有一、兩頁 empty 頁需要有浮水印 (如封面、書名頁)
%% 可以在該頁 (無頁碼) 使用 \thispagestyle{EmptyWaterMarkPage}
%% 等同於 empty page style + water mark
\fancypagestyle{EmptyWaterMarkPage}{%
\fancyhead{}%
\PlaceWaterMark%
\fancyfoot{}%
\renewcommand{\headrulewidth}{0pt}%
\renewcommand{\footrulewidth}{0pt}%
}

%% <<< 浮水印

% 如需額外的頁楣 (header) 或 footer，請在 headerfooter.tex 裡依例修改
% 它的預設內容是都關掉，可依需要打開
% this file is encoded in utf-8
% v3.0 (Jun. 11, 2019)

%%%%%%% 其他的 header (left, right) 定義
% 底下定義了一些常見的 header 型式
% 預設情況是關掉的
% 使用者可以視需要將之打開
% 也就是把下列介於 >>> 與 <<< 之間
% 的文字行打開 (行首去掉百分號)

%% header >>>
%\renewcommand{\chaptermark}[1]{%
%\markboth{\prechaptername\ \thechapter\ \postchaptername%
%\ #1}{}%
%}  %定義 header 使用的「章」層級的戳記
%\fancyhead[L]{} % 左 header 為空
%\fancyhead[R]{\leftmark}  % 右 header 擺放「章」層級的戳記 (以 \leftmark 叫出)
%\renewcommand{\headrulewidth}{0.4pt}  % header 的直線 0.4pt; 0pt 則無線
%% <<< header

%%%%%%% 其他的 footer (left, right) 定義
% 底下定義了一些常見的 footer 型式
% 預設情況是關掉的
% 使用者可以視需要將之打開
% 也就是把下列介於 >>> 與 <<< 之間
% 的文字行打開 (行首去掉百分號)

%% footer >>>
%\fancyfoot[L]{} % 左 footer 為空
%\fancyfoot[R]{\small{NTU \LaTeX\ v3.0}} % 右 footer 擺放論文格式版本
%\renewcommand{\footrulewidth}{0.4 pt} % footer 的直線 0.4pt; 0pt 則無線
%% <<< footer

%%%%%%%%%%%%%%%%%%%%%%%%%%%%%

% 常見的 notation 指令
% 請修改 notations.tex 加入自己需要的 notations
% %%%% Common Symbol %%%%
\newcommand*{\Nb}{\mathbb{N}}
\newcommand*{\Zb}{\mathbb{Z}}
\newcommand*{\Qb}{\mathbb{Q}}
\newcommand*{\Rb}{\mathbb{R}}
\newcommand*{\Cb}{\mathbb{C}}
\newcommand*{\Eb}{\mathbb{E}}

\newcommand*{\Ac}{\mathcal{A}}
\newcommand*{\Dc}{\mathcal{D}}
\newcommand*{\Oc}{\mathcal{O}}
\newcommand*{\Tc}{\mathcal{T}}
\newcommand*{\Uc}{\mathcal{U}}
\newcommand*{\Vc}{\mathcal{V}}
\newcommand*{\Xc}{\mathcal{X}}
\newcommand*{\Yc}{\mathcal{Y}}
\newcommand*{\Zc}{\mathcal{Z}}
\newcommand*{\trans}{\mathsf{T}}

%%%% mapping Symbol %%%%
\newcommand*\bij{\lhook\joinrel\twoheadrightarrow}
\newcommand*\oneto{\hookrightarrow}
\newcommand*\onto{\twoheadrightarrow}
\newcommand*\isoto{\xrightarrow{\sim}}
\newcommand*\acts{\curvearrowright}
\newcommand*\revacts{\curvearrowleft}

%%%% set definition %%%%
% just to make sure it exists
\providecommand\given{}
% can be useful to refer to this outside \Set
\newcommand*\SetSymbol[1][]{%
  \nonscript\:#1\vert
  \allowbreak
  \nonscript\:
\mathopen{}}
\DeclarePairedDelimiterX\Set[1]\{\}{%
  \renewcommand\given{\SetSymbol[\delimsize]}
  \,#1\,
}

%%%% Probability and Statistics Notations %%%%
\DeclarePairedDelimiterXPP{\KL}[2]{D_\text{KL}}\lbrack\rbrack{}{{#1} \delimsize\Vert {#2}}
\DeclarePairedDelimiterXPP{\Prob}[1]{\Pb}\lbrace\rbrace{}{#1}
\DeclarePairedDelimiterXPP{\Ev}[1]{\Eb}\lbrack\rbrack{}{#1}
\DeclarePairedDelimiterXPP{\Evr}[2]{\Eb_{#1}}\lbrack\rbrack{}{#2}

%%%%%%%%symbol and function settings%%%%%%%%%
\DeclarePairedDelimiter{\abs}{\lvert}{\rvert}
\DeclarePairedDelimiter{\norm}{\lVert}{\rVert}
\DeclarePairedDelimiter{\inpd}{\langle}{\rangle} % inner product
\DeclarePairedDelimiter{\ceil}{\lceil}{\rceil}
\DeclarePairedDelimiter{\floor}{\lfloor}{\rfloor}
\DeclareMathOperator*{\argmin}{arg\,min}
\DeclareMathOperator*{\argmax}{arg\,max}
\newcommand*{\bv}[1]{\mathbf{#1}}

%% some function name
\DeclareMathOperator{\ELBO}{ELBO}
\DeclareMathOperator{\ReLU}{ReLU}



% Title Page
% 請修改 title_settings.tex 設定論文題目/指導教授等變數
\renewcommand{\enTitle}{Research of Discrete Speech Representations}  %英文標題
\renewcommand{\zhTitle}{語音離散化表徵}  %中文標題
\renewcommand{\authorZhName}{陳建成}  %作者中文姓名
\renewcommand{\authorEnName}{Chien-Cheng Chen}  %作者英文姓名
\renewcommand{\authorStudentID}{R09942097}  %作者學號
\renewcommand{\advisorZhName}{李琳山}  %指導教授中文姓名
\renewcommand{\advisorEnName}{Lin-shan Lee}  %指導教授英文姓名
\renewcommand{\zhCollegeName}{電機資訊學院}  %學院中文名稱
\renewcommand{\enCollegeName}{College of Electrical Engineering and Computer Science}  %學院英文名稱
\renewcommand{\zhDepartmentName}{電信工程學研究所}  %系所中文名稱
\renewcommand{\enDepartmentName}{Graduate Institute of Communication Engineering}  %系所英文名稱
\renewcommand{\rocYear}{一百一十三}  %中華民國紀年年份
\renewcommand{\zhMonth}{六}  %中文月份
\renewcommand{\enYear}{2024}  %公元紀年
\renewcommand{\enMonth}{June}  %英文月份
\renewcommand{\oralDate}{113 年 6 月 1 日}  %口試日期


% 載入中文名詞的定義：例如，Figure -->「圖」, Chapter -->「第 x 章」
% this file is encoded in utf-8
% v3.0 (Jun. 11, 2019)

% 下列中文名詞的定義，如果以註解方式關閉取消，
% 則會以系統原先的預設值 (英文) 替代
% 名詞 \prechaptername 預設值為 Chapter
% 名詞 \postchaptername 預設值為空字串
% 名詞 \tablename 預設值為 Table
% 名詞 \figurename 預設值為 Figure
\renewcommand\prechaptername{第} % 出現在每一章的開頭的「第 x 章」
\renewcommand\postchaptername{章}
\renewcommand{\tablename}{表} % 在文章中 table caption 會以「表 x」表示
\renewcommand{\figurename}{圖} % 在文章中 figure caption 會以「圖 x」表示

% 下列中文名詞的定義，用於論文固定的各部分之命名 (出現於目錄與該頁標題)
\newcommand{\nameInnerCover}{書名頁}
\newcommand{\nameCommitteeForm}{論文口試委員審定書}
\newcommand{\nameCopyrightForm}{授權書}
\newcommand{\nameCabstract}{中文摘要}
\newcommand{\nameEabstract}{英文摘要}
\newcommand{\nameAckn}{誌謝}
\newcommand{\nameToc}{目錄}
\renewcommand{\contentsname}{目錄}
\newcommand{\nameLot}{表目錄}
\newcommand{\nameTof}{圖目錄}
\newcommand{\nameSlist}{符號說明}
\newcommand{\nameRef}{參考文獻}
\newcommand{\nameVita}{自傳}


% 如果不需要以中文數字一、二、三呈現章別，例如「第一章」
% 則請把下列介於 >>> 與 <<< 之間
% 的文字行關掉 (行首加上百分號), 會以「第 1 章」呈現
%% 中文數字章別 >>>
% this file is encoded in utf-8
% v3.0 (Jun. 11, 2019)

% 請依需要選擇其中一種表現方式，把它所對應的指令列打開，其他沒有用到的表現方式的對應指令列請關閉。(用行首百分號)

%% 第一種目錄格式：
%%	1  簡介 ............................ 1
%%
%%      章別 (chapter counter) 「1」前後沒有其他文字，
%%
%%      內文章標題是
%%		第 1 章	簡介
%%	\tocprechaptername, \tocpostchaptername 都設成沒有內容的空字串
%%	\tocChNumberWidth 設成 1.4em (預設)
%%      底下三行指令請打開
%\renewcommand\tocprechaptername{}
%\renewcommand\tocpostchaptername{}
%\setlength{\tocChNumberWidth}{1.4em}


%% 第二種目錄格式：
%%	一、簡介 ............................ 1
%%
%%      章別 (chapter counter) 「一」前沒有文字，後有頓號，
%%
%%      內文章標題是
%%		第一章		簡介
%%	\tocprechaptername 設成沒有內容的空字串
%%	\tocpostchaptername 設成頓號
%%	\tocChNumberWidth 設成 2em
%%      底下四行指令請打開 (預設)
\renewcommand\countermapping[1]{\zhnumber{#1}}
\renewcommand\tocprechaptername{}
\renewcommand\tocpostchaptername{、}
\setlength{\tocChNumberWidth}{2em}


%% 第三種目錄格式：
%%	第一章、簡介 ......................... 1
%%
%%      章別 (chapter counter) 「一」前有「第」，後有「章」與頓號，
%%      內文章標題是
%%		第一章		簡介
%%	\tocprechaptername 設成「第」
%%	\tocpostchaptername 設成「章、」
%%	\tocChNumberWidth 設成 3em
%%      底下四行指令請打開
%\renewcommand\countermapping[1]{\CJKnumber{#1}}
%\renewcommand\tocprechaptername{第}
%\renewcommand\tocpostchaptername{章、}
%\setlength{\tocChNumberWidth}{3em}



%% 可以依照需要作彈性的設定
%%
%% 章別 (數字，包括後面的字串) 的寬度 \tocChNumberWidth，
%% 會影響章名與章別之間的間隔 (太少則相疊，太多則留白)
%% 建議設成 \tocpostchaptername 內容字數加一，做為 em 的倍數，
%% 但至少也要有 1.4 倍。

%% <<< 中文數字章別

%%%%%%%%%%%%%%%%%%%%%%%%%%%%%
%  end of preamble
%%%%%%%%%%%%%%%%%%%%%%%%%%%%%
%
\begin{document}
\CJKindent  %%% ZZZ %%%  段首內縮兩格

%%% 以下載入前頁、本文、後頁

% \NTUtitlepage  % 產生論文封面

\newpage
\setcounter{page}{1}
\pagenumbering{roman}

\NTUoralpage  % 產生口試委員會審定書

\mydoublespacing
\begin{acknowledgement} %誌謝
    兩年來的碩士生活，隨著這份論文的誕生而告了一個段落。回思過去兩年來的生活，有做出成果的喜悅、也有處處碰壁時的苦悶、也有跟同學共同奮鬥的記憶，自己在這兩年中實在成長了許多，而這些都要感謝實驗室的大家長：李琳山教授。教授在實驗室營造了自由研究的氛圍，讓實驗室的同學都能按自己喜好自由發展研究方向，並從旁關心協助同學的研究。我在這樣的氛圍下也受益許多，學習到了許多做研究與做人處事的方法。

    這些研究能夠順利完成要由衷地感謝我的家人，他們無論何時都支持我的決定，並且從旁給我協助，在我最忙碌而都很晚回家的那段時間，他們也是很包容我，並給予我支持。

    與實驗室的同學相處的這段時光將是碩士生活中最難忘的一段日子，我們無論是修課或是研究上都是彼此的好戰友，我從你們身上都學到很多；下一屆的你們將是實驗室下一代的主力，祝你們未來研究順利！

    最後要感謝我的朋友、同學、以及伙伴們，不論是平常一起吃飯聊天、有正事時的一起奮鬥、或一起出去玩，你們都是我平常生活上最大的支持！

\end{acknowledgement}

\begin{zhAbstract} % 中文摘要
    這是中文摘要
\end{zhAbstract}

\begin{enAbstract} % 英文摘要
    This is an English abstract.
\end{enAbstract}

{
%\zhKaiFont
\mysinglespacing\selectfont
\tableofcontents %目錄

\listoffigures  %圖目錄

\listoftables  %表目錄
\par
}

\newpage
\setcounter{page}{1}
\pagenumbering{arabic}
  % omitted as...: 
% <<<
% \NTUtitlepage  % 產生論文封面

% \newpage
\setcounter{page}{1}
\pagenumbering{roman}

% \NTUoralpage  % 產生口試委員會審定書

\mydoublespacing
%% \begin{acknowledgement} %誌謝
%%     我不想看到

%% \end{acknowledgement}%% 

%% \begin{zhAbstract} % 中文摘要
%%     這是中文摘要
%% \end{zhAbstract}%% 

%% \begin{enAbstract} % 英文摘要
%%     This is an English abstract.
%% \end{enAbstract}

%%%%%% 先藏起來
%%%%%%%% {
%%%%%%%% %\zhKaiFont
%%%%%%%% \mysinglespacing\selectfont
%%%%%%%% \tableofcontents %目錄
%%%%%%%% 
%%%%%%%% % \listoffigures  %圖目錄
%%%%%%%% 
%%%%%%%% % \listoftables  %表目錄
%%%%%%%% \par
%%%%%%%% }
%%%%%% <<<

\newpage
\setcounter{page}{1}
\pagenumbering{arabic}

% >>>
% % !TeX root = ../thesis.tex

% !TeX root = ../../thesis.tex

\newcommand{\NOTEHERE}[1]{\colorbox{yellow}{\textcolor{blue}{
    \lparen\lparen\lparen}\textcolor{
        red}{#1}\textcolor{blue}{\rparen\rparen\rparen}}}

% newline
\newcommand{\PROBABLY}[1]{}
%% \newcommand{\PROBABLY}[1]{\colorbox{red}{\textcolor{green}{
%%     \lbrack\lbrack\lbrack}\textcolor{green}{#1}\textcolor{
%%         green}{\rbrack\rbrack\rbrack}}}

% \textcolor{green}{#1}
\newcommand{\FIXTHIS}[1]{\par\makebox[\textwidth]{
    \colorbox{pink}{\textcolor{green}{
        \lbrack\lbrack\lbrack}\textcolor{
            blue}{#1}\textcolor{green}{\rbrack\rbrack\rbrack}}}}

\newcommand{\mysetcounter}[2]{\setcounter{#1}{\the\numexpr#2 - 1\relax}}
\newcommand{\mychcnt}[1]{\mysetcounter{chapter}{#1}}

% % !TeX root = ../thesis.tex
\chapter{導論}
\section{研究動機}

　　語言是人與人彼此交流最主要的橋樑，而人們互相溝通最自然的方式便是透過說話的語音（speech）達成。人類往往是自幼就牙牙學語開始說話，直到已屆學齡左右才開始學習認字與書寫。雖然在這個資訊爆炸的時代，人們已經習慣以文字呈現的語言作為獲取資訊的主要媒介，但不論如何，各種書寫系統其背後承載的語言必定有語音的形式作為對應。更何況世界上現存大約七千多種  \cite{ethnologue}  語言中，絕大多數不見得存在成熟且普及的文字系統，卻無礙於這些語言被人們所熟悉和使用。因此，「語音」作為語言不可或缺的存在方式，了解它和研究它的價值自然不言而喻。

　　然而，相對於穩定、易於處理和保存的文字文本，語音訊號的變化萬千，蘊藏了大量從語者風格、表達內容到抑揚頓挫（韻律，prosody）等不同層次的訊息，使得對它的處理、研究相比之下複雜度與難度劇增。由於語音的這種特性，過往對於語言最有興趣的語言學家們，即便明白語音作為多數語言主體的事實，也不得不藉文字符號為依託進行探索。進入資訊化時代後，藉助電腦硬體等計算設備的幫助，從語料庫、計算語言學到自然語言處理等透過科技的力量發展語言處理技術的領域，頗長一段時間也是專注於文字的處理與分析。
而嘗試結合訊號處理發展的語音技術領域，當時則是透過語言學家對語言的領域知識，例如從 phoneme、morphology、語法等等用以刻劃人類語音和語言特性的概念，將之結合機器學習建立模型，開發技術以方便人們能以語音這種更靈活的媒介，更好的讓電腦、手機等科技工具可以更接近「直接溝通」的使用方式，便利人們的日常生活。

　　近年來，由於圖形處理器（graphics processing unit，GPU）等硬體平行運算技術的進步，深層學習（deep learning）快速崛起成為人工智慧的主流，有了此項機器學習的技術，模型的彈性能夠更好的萃取資料、更貼近的尋找資料背後的機制並進行預測，使得人們不再非得依賴大量費時費工的人類標注過程，進而使得利用大量語料庫發展語言技術，進一步推進語言科技發展成為可能。尤其在自監督學習（self-supervised learning）技術出現之後，深層學習模型可以依照人們給定的方向，更細緻的從大量未標注、相較容易取得的語音或文字的語料，找出其中的語音、語法及語義等等結構，形成帶有對人類語言有前所未見表現的基石模型（foundation model），是這個領域的一大里程碑。尤其在以處理文字為主體的自然語言處理領域，甚至出現了幾乎使人類真偽難辨的生成式模型，改變了人們生活的方方面面。

＿＿借鏡文字方面的成功經驗，語音處理領域的研究者們也開始嘗試將語言模型（language model）的概念套用於變化莫測的語音訊號之上，原先人們藉助訊號處理知識一直使用的各種語音訊號特徵（feature）也在自監督學習的架構之下，出現了許多模型從大量語音資料中得到的「語音表徵（speech representation）」，作為精煉語音資訊的另外一種新選擇開始廣泛被採用。然而，相比於文字符號的穩定與單純，語音的複雜性使得它處理起來會需要更大量的資料和運算資源來擷取其中不同層次的細節，而且作為物理訊號，語音還必須處理掉環境中的雜訊等干擾。為了從紛亂的聲音中提取出最重要的訊息，向量量化（vector quantization）的技巧因而經常被使用在語音 \cite{chorowski2019unsupervised, chen2023vector, zhao2023speech} 或影像的領域中。爾後， \cite{lakhotia2021generative}  基於模仿人類學習語言的過程，藉助諸如 CPC（\cite{oord2019representation} ）、HuBERT （\cite{hsu2021hubert} ）、wav2vec 2.0 （\cite{baevski2020wav2vec} ）等自監督學習模型的幫助，引入向量量化的技術，提出了「無文字（textless）」的學習架構，轉而以語音表徵量化後的「離散單元（discrete unit）」作為操作對象，企圖以單純大量的語音資料中訓練出一個不依賴文字的語言模型。此種學習架構的優勢在於在能保有利用大量未標注文字轉寫語音資料的同時，與連續表徵相比資訊的位元率（bit rate）利用更有效率、容易儲存、處理與傳輸，以及形式上更像文字的特性。自此之後有一系列  \NOTEHERE{補足 unit 相關的續作，例如英翻臺} 使用 discrete units 進行任務訓練的研究，一定程度的印證了unit 捕捉語音內容的有效性。

　　Unit 固然在編碼語音之上有不錯的效果，因此 empirically 被領域內的不少學者直接視為「機器自己學習出來的文字」，可是究竟這項 claim 本身有多符合實際狀況，有學者便嘗試重新從語言學的概念汲取靈感，對其進行語音學層面的分析析析析 \NOTEHERE{放唐顥等等的語音學研究}你知道吧？
　　\PROBABLY{你可能要併入連續的 speech repr}
　　\FIXTHIS{確認一下 acoustic piece 那邊的文線是不是也是如我那樣說的}
　　\FIXTHIS{確認研究動機要寫到哪裡}


\section{研究方向}

　　本論文旨在探討自監督學習模型的離散語音表徵與音位（phoneme）之間的關係。

\section{主要貢獻}
\section{章節安排}

本論文之章節安排如下：

% \chapter{導論}
\section{研究動機}

　　語言是人與人之間最主要的溝通方式，是我們對外界世界互動最主要的媒介。隨著時代的發展，在這個資訊爆炸的社會，人們接觸到語言、使用與互動的情形變得非常頻繁。然而，語言的多樣性在此卻成為一項阻礙不同人們交流的障礙，因此，發展語言科技為人們的互動提供協助，因而成為了一項不可避免的趨勢。藉助科技與知識的力量，人們可以透過不同語言的視角去獲取更多元的資訊，激發思考、開拓視野，因此諸如語音辨識、機器翻譯等等，一直都是人們熱衷於研究的議題。

    過往，為了達成此一遠大目標，機器學習工程師們與語言學家之間互相合作，期望藉由對領域知識的了解去建立模型，打造出創新的技術以應對語言的各種變化，來滿足人們溝通的需求。近年來，由於硬體平行運算技術的進步，深層學習（deep learning）快速崛起成為人工智慧的主流，有了此項機器學習的技術，模型的彈性能夠更好的萃取資料、更貼近的尋找資料背後的機制並進行預測，使得人們不再非得依賴大量費時費工的人類標註過程，進而使得利用大量語料庫發展語言技術，進一步推進語言科技發展成為可能。尤其在自監督學習（self-supervision）技術出現之後，深層學習模型可以依照人們給定的方向，更細緻的從大量未標注、相較容易取得的語音或文字的語料，找出其中的語音、語法及語義等等結構，形成帶有對人類語言有前所未見表現的基礎模型（foundation model），是這個領域的一大里程碑。尤其在以處理文字為主體的自然語言處理領域，甚至出現了幾乎使人類真偽難辨的生成式模型，改變了人們生活的方方面面。

    ……………………………………………………………………………… \\
    ………………………………………………………………………………………… \\
　　………………………………………………………………………………

    然而相較於穩定、易於處理的文字文本，語音訊號的變化複雜萬千，蘊藏了大量如內容、韻律、語者等等不同層次的訊息，使得處理的難度劇增。更何況，目前世界上大約七千多種語言中，絕大多數仍然沒有成熟且普及的文字系統。

    （先從語音表徵開始）

    因此，除了發展文字的模型以外，語音處理的技術更是必不可少的。
而在語音處理這邊，近期藉由自監督學習，提出了許多語音這邊的基礎模型。
% 這部分跟前面重疊到了
雖然

    因此在語音處理界，「無文字（textless）」的發展是相當吸引人的。

    （加上語言公平性問題？）

    近期由於 HuBERT 等等起來的 GSLM 等架構，已經在一定程度上做到完全不依賴文字轉寫、單純在大量的語音語料之上建立一定程度上媲美於「大型語言模型」的成果，甚至達成了閩南語和英語的互相對譯。此一里程碑大力的推動了語言科技的進展，有望推動藉助科技達成的降低語言障礙。

    ………………………………………………………………………………

    
    
   不過，即便語音的技術已經相當成熟，在追求模型表現的同時，語音與語言的技術開始與過往的人類對於語言學、語音的理解逐漸產生脫節。似乎在為了讓機器可以擁有良好表現的同時，理解模型如何運作似乎是被犧牲的必然。但人們在追求更好的模型表現的同時，有一群人開始注意到語音處理模型是否有可能抓取到人類語音中特有的、區別於一般音訊的特徵，並嘗試使用過往用來研究、歸類人類語音的方式，結合機器學習與統計學去解釋為什麼，並期求可以比較甚至改善機器模型在進行語音處理時的表現，不僅僅只是使用資料集本身的分數，而有更多更多元、更穩健的衡量標準。由於離散表徵在當今語音模型與語音處理技術已經愈來愈具備重要性，因此探討與分析為什麼語音離散表徵可以幫助下游任務的背後成因是相當重要的研究方向，其中一個驗證離散表徵能夠幫助模型處理語音訊號的方式，便是驗證其與音位（phoneme）之間的對應關係。

    我們想要知道的是，在離散單元推動語音處理發展的同時，它究竟與人類書寫和使用的文字還有多遠的差異，以及在使用上是否能夠達成如同文字的效果，仍舊是領域中尚待探討的議題。所以我們先看看，他是不是起碼符合語音學上，作為「捕捉語音內容」的基本特徵。

    （加上跟 LLM 的關係嗎？）

    也就是Phonology去討論這件事，看它跟聲學特徵跟語音到底有多像，從這個方向我們可以進而去改善說，也許我們可以去用不同的方式使用Unit，來進一步推進讓我們的Unwritten Language這些技術，能夠更好的跟，就是等於說進而發揮類似文字的效果，然後同時也能促進我們去知道，機器是怎麼樣去理解這些語音訊號，它可能看的是哪些部分。

    ……………………………………………………………………………… \\
    ………………………………………………………………………………………… \\
　　………………………………………………………………………………


\section{研究方向}

%%%     於是我們對這些離散單元進行了單一的分析
%%%     另外，因為人類的音素往往是超過一個音框的單位
%%%     因此推測離散單元單一本身會抓取到的是更加細緻的特徵
%%% 
%%%     那為了找出更加符合人類文字認知的表示符號
%%%     因此我們採用文字那邊常用的分詞器
%%%     // 可能要紀錄 rate 來證明這件事了
%%%     嘗試看看可否找出更接近人們認知信息率的符號，以其這種機器抓出
%%% 來的 fine structure 可以更往人類的文字靠近
%%% 
%%%     具體方法是，針對這些抓出來的 unit 跟 piece, 我們去算 purtiy,
%%% pmni, segmentation 這些指標，看看他們是不是真的跟人類 top down
%%% 標註出來的的 phoneme label 有統計上的一致性。
%%% 
%%%     最後我們嘗試做在 PR 或 ASR 上，驗證這個一致性不是僅僅在統計數
%%% 據上顯現，也是可以在應用上得到幫助的。

　　在近年，已經有不少相關的研究開始嘗試往將離散單元（discrete unit）作為除了連續表徵（continuous representation）以外，可以編碼（encode）語音訊號的另外一種方式。離散表徵（discrete representation）跟連續表徵相比，具有資訊更濃縮（位元率（bit rate）更低）因此更好儲存、處理與傳輸，以及形式上更像文字的特性。

    儘管離散表徵在語音社群（community）中常被當成一種類似文字的存在，另外有一些文獻則是將其當成連續表徵的精簡表示法。

    然而

    因此，我們藉由分析各種離散單元和人類理解語音最直接的處理層次：音位（phoneme）之間的關係，並將兩者進行各種統計上的序列比較，可以作為訓練大型語音基石模型（foundation model）的分詞（tokenization）基礎，選定最適合的表徵最為系統的輸入符記（token）。

    如 \cite{wu2023wav2seq} 等作品便是首先嘗試將離散單元進行 tokenization 後做進一步處理的，其後的


\section{主要貢獻}  % 有什麼

　　結果我們發現，藉由觀察這些 unit 並嘗試藉由分詞演算法找出更 high-level 的單位之後，我們觀察到這些機器學習 figure out 的「偽標記」一定程度上的符合了人類音素的特性，因此可以當成某種類似拼音文字的存在。當然這跟人類真正使用的拼音文字仍有距離，
% // 難道可以探討 text 是不是一種 speech tokenization 過程嗎？
因此雖然無法直接當成 exact 的文字使用，一些人類的標注還是需要的，不過透過 ML 我們已經可以盡量更有效的利用珍貴的標注資料，幫助那些尚不容易取得文字的語言發展語音語言科技，以協助保存他們的語言。
% // 問題是，有些語音任務並不需要抽內容嗎？


\section{章節安排}

　　由於本論文是以剖析既有的語音離散表徵為主軸，因此就相關研究方面需要從各角度入手，單獨成一章節。接著我們會從單一的離散單元，以及將單元視為像文字的字符（character）並進行分詞演算法兩種對語音離散單元處理的層次分別成章進行分析，最後將這些表徵嘗試做在語音的任務上，以驗證其具有一定的語音表徵能力，且能保留語音學的特徵。



%%%%%%%%%%%%%%%%%%%%%%%%%%%%%%%%%%%%%%%%%%%%%%%%%%%%%%%%%%%%%%%%%%%%%
%%%%%%%%%%%%%%%%%%%%%%%%%%%%%%%%%%%%%%%%%%%%%%%%%%%%%%%%%%%%%%%%%%%%%
%%%%%%%%%%%%%%%%%%%%%%%%%%%%%%%%%%%%%%%%%%%%%%%%%%%%%%%%%%%%%%%%%%%%%

\chapter{背景知識}

\section{深層類神經網路（deep neural network）}

深層類神經網路（Deep neural network） 是麥氏（McCulloch）在 1943 年提出 \cite{mcculloch1943logical} 仿生數學模型，旨在模擬生物神經系統的連結。

深層類神經網路是一個取法自生物神經連結的數學模型，其在計算認知神經科學中以連結派（connectionism）為主要代表，後在電腦科學與機器學習中有不同結構的進展。在此之後，基於其彈性與平行化的能力，能在 GPU 上面很有效率的進行運算並達到前所未有的效能，因此現在已經成為人工智慧發展的主流。

基於深層類神經網路的神經架構有 CNN、RNN、Transformer 等等，由於這些架構在語音與文字處理上都已經被廣泛使用，因此在下面分別介紹：


\subsection{卷積式（convolutional）類神經網路}

卷積式類神經網路一開始是在 \\cite{} 中提出，主要是鑑於影像中的局部性（locality），讓 NN 可以在

。

在語音中，因為

語音訊號的資訊是被呈現在時間維度上
，
因此
通常使用

一維

的
卷積式類神經網路
，
以捕捉
時間維度上的局部性特徵
，例如本研究特別探討的 phoneme、morpheme 等等。

% \subsection{遞迴式（recurrent）類神經網路}



\subsection{
    序列至序列（sequence-to-sequence）模型}

    


\subsection{專注（attention）機制}    

\subsection{轉換器（Transformer）}



\section{表徵（representation）學習}


\subsection{文字的語意表徵}

\subsection{語音特徵與表徵}

\section{語音基石模型與自監督式學習}

\subsection{自監督式學習}

\subsection{語音基石模型}

\subsection{離散單元}

\section{本章總結}

\setcounter{chapter}{1}
<<<<<<< HEAD
% !TeX root = ../thesis.tex
\mychcnt{3}\chapter{
單一語音離散表徵與語音標記的對應模式}
\section{動機}  % \section{簡介}
\section{相關研究}

在 HuBERT 出來之後，

\subsection{語音表徵的語音學分析}
\section{衡量方式}
\subsection{音位（phoneme）長條圖（bar chart）}
\subsection{純度（purity）、熵（entropy）}
\subsection{對齊（alignment）}
\section{語音學分類}
\section{分析結果}
\subsection{基於各自音位的分析}
\subsection{基於語音學分類的分析}
\section{本章總結}

% ============================================================ %
% ============================================================ %

\mychcnt{2}\chapter{
背景知識}
\section{深層類神經網路（deep neural network）}

　　深層類神經網路（Deep neural network） 是麥氏（McCulloch）在 1943 年提出 \cite{mcculloch1943logical}，取法自生物神經連結的計算模型，旨在模擬生物神經系統的連結，以模仿生物的各項功能，進而透過機器學習的最佳化演算法，使得整個模型能夠藉由資料去貼合理想的函數，以達成應用或工程上所需要的各種任務。 。

    以發展此模型為主軸的心理學流派，在計算認知神經科學中被稱為「連結派（connectionism）」，其後因為該網路的彈性與平行化的能力，和諸如圖形處理器（graph processing unit，GPU）等硬體裝置能夠 最好的利用。 
    并能夠更 好的描述資料分佈、達到前所未有的效能 ，因此近年在電腦科學的機器學習領域中獲得重大進展，並因此現已成為人工智慧發展的主流。

    基於深層類神經網路的神經架構有 CNN、RNN、Transformer 等等，由於這些架構在語音與文字處理上都已經被廣泛使用，因此在下面分別介紹：

\subsection{卷積式（convolutional）類神經網路}

　　卷積式類神經網路一開始是在 cite 中提出，主要是鑑於影像中的局部性（locality），讓 NN 可以在。在語音中，因為語音訊號的資訊是被呈現在時間維度上，因此通常使用一維的卷積式類神經網路，以捕捉時間維度上的局部性特徵，例如本研究特別探討的 phoneme、morpheme 等等。

\subsection{遞迴式（recurrent）類神經網路}

\subsection{
    序列至序列（sequence-to-sequence）模型}
    
    由于許多實際上的資料都是2個序 之間互相配對的關系 此類的資料包含語音文字。 信號等等，都是以時間軸為主要 演變方向的資料。 因此 有一類模型。 會被以序列制序列的模式進行訓練。 旨在模擬輸入與輸出序列之間的變化與相依關係（dependency）。

此類模型一般的架構是由一個編碼器和解碼器構成 其中編碼器是將輸入訊號借由內部表征進行編碼。 依據每個時間點輸入訊號的順序來改變其內部表征的狀態 接著將最後一個時間點的表徵作為整個序列的特征傳遞給解碼器進行輸出訊號 生成。 
    
\subsection{專注（attention）機制}
    
    原本的序列自序列模型本身。 需要讓解碼器單純透過最後一個時間 點。 的表征資訊來完全儲存輸入序列的一切資訊 以工解碼器判斷。 并生成輸出序列。 然而，由于 單就最後一個向量進行判斷對於解碼器而言過於不易。因此。 盧氏提出。 在編碼器中對輸入序列的不同時間點進行注意力機制 亦即讓解碼器可以根據當下 所需要輸出的 內容判斷應該要重新對輸入序列的哪些部份 進行更多的加權 。

    

\subsection{轉換器（Transformer）}

其後，由瓦氏（Vaswani） cite 提出的 論文中 提出了一個完全由注意機制。 所構成的序列自序列模型。 原先該模型適用於解決機器翻譯。 的問題。
由於其能夠高度平行化的特性， 日後在 自然源處理和語音處理，甚至到電腦視覺領域等近乎整個深層學習的領域都被廣泛的應用。  


\section{表徵（representation）學習}

\subsection{文字的語意表徵}

word2vec 是 mikolov \CITEME 最早提出跟對於文字進行語義表徵的 work。在其後

\subsection{語音特徵與表徵}


\section{語音基石模型與自監督式學習}

\subsection{自監督式學習}

\subsection{語音基石模型}

\subsection{離散單元}


\section{本章總結}



% % !TeX root = ../thesis.tex
\mychcnt{3}\chapter{
    單一語音離散表徵與語音標記的對應模式}
\section{動機}  % \section{簡介}
\section{相關研究}

在 HuBERT 出來之後，

\subsection{語音表徵的語音學分析}
\section{衡量方式}
\subsection{音位（phoneme）長條圖（bar chart）}
\subsection{純度（purity）、熵（entropy）}
\subsection{對齊（alignment）}
\section{語音學分類}
\section{分析結果}
\subsection{基於各自音位的分析}
\subsection{基於語音學分類的分析}
\section{本章總結}

\mychcnt{2}\chapter{背景知識}
\section{真全部深層類神經網路（deep neural network）}

　　深層類神經網路（Deep neural network） 是麥氏（McCulloch）在 1943 年提出 \cite{mcculloch1943logical}，取法自生物神經連結的計算模型，旨在模擬生物神經系統的連結，以模仿生物的各項功能，進而透過機器學習的最佳化演算法，使得整個模型能夠藉由資料去貼合理想的函數，以達成應用或工程上所需要的各種任務。 。

    以發展此模型為主軸的心理學流派，在計算認知神經科學中被稱為「連結派（connectionism）」，其後因為該網路的彈性與平行化的能力，和諸如圖形處理器（graph processing unit，GPU）等硬體裝置能夠 最好的利用。 
    并能夠更 好的描述資料分佈、達到前所未有的效能 ，因此近年在電腦科學的機器學習領域中獲得重大進展，並因此現已成為人工智慧發展的主流。

    基於深層類神經網路的神經架構有 CNN、RNN、Transformer 等等，由於這些架構在語音與文字處理上都已經被廣泛使用，因此在下面分別介紹：

\subsection{卷積式（convolutional）類神經網路}

　　卷積式類神經網路一開始是在 cite 中提出，主要是鑑於影像中的局部性（locality），讓 NN 可以在。在語音中，因為語音訊號的資訊是被呈現在時間維度上，因此通常使用一維的卷積式類神經網路，以捕捉時間維度上的局部性特徵，例如本研究特別探討的 phoneme、morpheme 等等。

\subsection{遞迴式（recurrent）類神經網路}

\subsection{
    序列至序列（sequence-to-sequence）模型}
    
    由于許多實際上的資料都是2個序 之間互相配對的關系 此類的資料包含語音文字。 信號等等，都是以時間軸為主要 演變方向的資料。 因此 有一類模型。 會被以序列制序列的模式進行訓練。 旨在模擬輸入與輸出序列之間的變化與相依關係（dependency）。

此類模型一般的架構是由一個編碼器和解碼器構成 其中編碼器是將輸入訊號借由內部表征進行編碼。 依據每個時間點輸入訊號的順序來改變其內部表征的狀態 接著將最後一個時間點的表徵作為整個序列的特征傳遞給解碼器進行輸出訊號 生成。 
    
\subsection{專注（attention）機制}
    
    原本的序列自序列模型本身。 需要讓解碼器單純透過最後一個時間 點。 的表征資訊來完全儲存輸入序列的一切資訊 以工解碼器判斷。 并生成輸出序列。 然而，由于 單就最後一個向量進行判斷對於解碼器而言過於不易。因此。 盧氏提出。 在編碼器中對輸入序列的不同時間點進行注意力機制 亦即讓解碼器可以根據當下 所需要輸出的 內容判斷應該要重新對輸入序列的哪些部份 進行更多的加權 。

    

\subsection{轉換器（Transformer）}

其後，由瓦氏（Vaswani） cite 提出的 論文中 提出了一個完全由注意機制。 所構成的序列自序列模型。 原先該模型適用於解決機器翻譯。 的問題。
由於其能夠高度平行化的特性， 日後在 自然源處理和語音處理，甚至到電腦視覺領域等近乎整個深層學習的領域都被廣泛的應用。  


\section{表徵（representation）學習}

\subsection{文字的語意表徵}



\subsection{語音特徵與表徵}


\section{語音基石模型與自監督式學習}

\subsection{自監督式學習}

\subsection{語音基石模型}

\subsection{離散單元}


\section{本章總結}



% % !TeX root = ../thesis.tex
\mychcnt{3}\chapter{
單一語音離散表徵與語音標記的對應模式}
\section{動機}  % \section{簡介}
\section{相關研究}

在 HuBERT 出來之後，

\subsection{語音表徵的語音學分析}
\section{衡量方式}
\subsection{音位（phoneme）長條圖（bar chart）}
\subsection{純度（purity）、熵（entropy）}
\subsection{對齊（alignment）}
\section{語音學分類}
\section{分析結果}
\subsection{基於各自音位的分析}
\subsection{基於語音學分類的分析}
\section{本章總結}

% ============================================================ %
% ============================================================ %

\mychcnt{2}\chapter{
背景知識}
\section{深層類神經網路（deep neural network）}

　　深層類神經網路（Deep neural network） 是麥氏（McCulloch）在 1943 年提出 \cite{mcculloch1943logical}，取法自生物神經連結的計算模型，旨在模擬生物神經系統的連結，以模仿生物的各項功能，進而透過機器學習的最佳化演算法，使得整個模型能夠藉由資料去貼合理想的函數，以達成應用或工程上所需要的各種任務。 。

    以發展此模型為主軸的心理學流派，在計算認知神經科學中被稱為「連結派（connectionism）」，其後因為該網路的彈性與平行化的能力，和諸如圖形處理器（graph processing unit，GPU）等硬體裝置能夠 最好的利用。 
    并能夠更 好的描述資料分佈、達到前所未有的效能 ，因此近年在電腦科學的機器學習領域中獲得重大進展，並因此現已成為人工智慧發展的主流。

    基於深層類神經網路的神經架構有 CNN、RNN、Transformer 等等，由於這些架構在語音與文字處理上都已經被廣泛使用，因此在下面分別介紹：

\subsection{卷積式（convolutional）類神經網路}

　　卷積式類神經網路一開始是在 cite 中提出，主要是鑑於影像中的局部性（locality），讓 NN 可以在。在語音中，因為語音訊號的資訊是被呈現在時間維度上，因此通常使用一維的卷積式類神經網路，以捕捉時間維度上的局部性特徵，例如本研究特別探討的 phoneme、morpheme 等等。

\subsection{遞迴式（recurrent）類神經網路}

\subsection{
    序列至序列（sequence-to-sequence）模型}
    
    由于許多實際上的資料都是2個序 之間互相配對的關系 此類的資料包含語音文字。 信號等等，都是以時間軸為主要 演變方向的資料。 因此 有一類模型。 會被以序列制序列的模式進行訓練。 旨在模擬輸入與輸出序列之間的變化與相依關係（dependency）。

此類模型一般的架構是由一個編碼器和解碼器構成 其中編碼器是將輸入訊號借由內部表征進行編碼。 依據每個時間點輸入訊號的順序來改變其內部表征的狀態 接著將最後一個時間點的表徵作為整個序列的特征傳遞給解碼器進行輸出訊號 生成。 
    
\subsection{專注（attention）機制}
    
    原本的序列自序列模型本身。 需要讓解碼器單純透過最後一個時間 點。 的表征資訊來完全儲存輸入序列的一切資訊 以工解碼器判斷。 并生成輸出序列。 然而，由于 單就最後一個向量進行判斷對於解碼器而言過於不易。因此。 盧氏提出。 在編碼器中對輸入序列的不同時間點進行注意力機制 亦即讓解碼器可以根據當下 所需要輸出的 內容判斷應該要重新對輸入序列的哪些部份 進行更多的加權 。

    

\subsection{轉換器（Transformer）}

其後，由瓦氏（Vaswani） cite 提出的 論文中 提出了一個完全由注意機制。 所構成的序列自序列模型。 原先該模型適用於解決機器翻譯。 的問題。
由於其能夠高度平行化的特性， 日後在 自然源處理和語音處理，甚至到電腦視覺領域等近乎整個深層學習的領域都被廣泛的應用。  


\section{表徵（representation）學習}

\subsection{文字的語意表徵}

word2vec 是 mikolov \CITEME 最早提出跟對於文字進行語義表徵的 work。在其後

\subsection{語音特徵與表徵}


\section{語音基石模型與自監督式學習}

\subsection{自監督式學習}

\subsection{語音基石模型}

\subsection{離散單元}


\section{本章總結}



% % !TeX root = ../../thesis.tex

% % !TeX root = ../thesis.tex

\chapter{
單一語音離散表徵與語音標記的對應模式}
\section{動機簡介}  % \section{簡介}
由於 HuBERT 之後，unit 的使用很廣泛，因此為了研究 unit 本身為什麼會被如此適當的可以讓模型視為文字對語音資料進行訓練，我們先從離散表徵本身的特徵分析起。 

% // 從為了驗證它是不是文字，改成分析內部有沒有精細結構
% ?? 要不要去分析有沒有 dseduplicate 這件事對於影響精細結構的關係，以及 rate 等等的差異呢？

\section{相關研究}
近期已經有多項相關的研究，嘗試在 SSL 這麼厲害的表現之後找原因，因此有針對 unit 背後 repr 的特性進行分析的 work，例如 CITEUSPLEASE。 

\subsection{語音表徵的語音學分析}
在 HuBERT 出來之後，有一些研究像是 cite 等等，試圖探討對於語音表徵這樣語音模型的基礎進行各種從統計和語音學領域知識角度的分析，以期望能夠解釋為什麼模型可以擁有如此的表現。

此後，cite 等等作品則是從原先連續的表徵出發，開始往離散的量化向量，甚至是離散單元進行分析比對。雖然分析的切入角度可以相當多樣，例如 ABX、tsne 降維分群等等，但本次研究主要著重比對兩者之間在同一段語音序列上給予標籤的相關性，也就是以「偽標籤（pseudo-label）」的角度進行衡量。

\subsection{無文字（textless）語言模型}

這系列 textless 以 GSLM 為最主要代表作，旨在探討 unit 作為一種替代文字的方案。

本論文以 GSLM textless 採用的模型 units 為主要分析對象，企圖銜接兩者的脈絡，來佐證這些 unit 作為一種「類似或可替代文字的語音紀錄方式」在能夠發揮 LM 的特長背後，是否是基於符合語音學特徵帶來的，抑或有什麼其他特徵。

\section{衡量方式（metric）}

為了測量這些 unit 跟 phn 這類語言學 labels 之間的 correlation，我們需要先介紹本論文會探討的指標


\subsection{音位（phoneme）長條圖（bar chart）}
要先對 unit 做頻率上的統計
\subsection{純度（purity）}
phoneme purity: 每個cluster unit內的phoneme purity，代表此unit是否有phoneme代表性
cluster purity：每個phoneme對應到的cluster統計，若cluster purity低代表less liguistical meaning?（抄hubert paper）單一 phoneme 本身對應的 unit 的一致性。如果

\subsection{熵（entropy）和相互資訊（mutual information，MI）}
除了「最大」的對應關係，根據 Hubert 原先的 paper CITEME (hubert) 中的分析方式，我們也可以從 info theory 的角度，去探討「觀察到一個 unit 對於 label 不確定性的降低」來考慮 unit 本身提供了多少背後 phn 的資訊

\subsection{對齊（alignment）}
cluster是否保留segment資訊，不將不同phoneme合併
segmentation → 
怕說只是每個 frame 本身 unit 或 piece 可以跟 phoeneme 特徵相關，但放在連續的語音中被切得很碎/前後文相關的東西不知道有沒有抓到


\section{語音學分類（phone type）}

\subsection{簡介}

除了單一 phn 本身的特性以外，由於 phn 本身彼此不是完全獨立的，而是彼此之間就存在相似的特徵，可以分成幾個組別。因此，依照 CITEME (tanghao等三篇) 的分組方式，對英語的 phn 進行分類並合併比對數據，看看這些 unit 本身是否有 capture 到相似的發聲特徵，而不單純只是把 phn 分成約五十類完全獨立的標籤。（基於語音表徵本身就是 acoustic sisgnals 來的，應該 by nature 要可以對語音特徵分組吧？）

以下為各分組進行簡單介紹：

\subsubsection{consonants}

子音可以分成五類

\subsubsection{vowels}

母音在這邊為了簡單起見，會被分在一起？

\subsection{解釋意義}

\begin{itemize}
    \item 純度（purity）：換成type之後有何變化（關聯性更強？）
    \item 熵（entropy）（放直方圖解釋） --> phone type更明顯？
    \item 對齊（alignment）：是否減少segment資訊的保留（連續子音母音被合併？
\end{itemize}

\section{分析結果}

 
\subsection{基於各自音位的分析}
\subsection{基於語音學分類的分析}
\section{本章總結}

=======
% !TeX root = ../../thesis.tex

% !TeX root = ../thesis.tex

\chapter{
單一語音離散表徵與語音標記的對應模式}
\section{動機簡介}  % \section{簡介}
由於 HuBERT 之後，unit 的使用很廣泛，因此為了研究 unit 本身為什麼會被如此適當的可以讓模型視為文字對語音資料進行訓練，我們先從離散表徵本身的特徵分析起。 

% // 從為了驗證它是不是文字，改成分析內部有沒有精細結構
% ?? 要不要去分析有沒有 dseduplicate 這件事對於影響精細結構的關係，以及 rate 等等的差異呢？

\section{相關研究}
近期已經有多項相關的研究，嘗試在 SSL 這麼厲害的表現之後找原因，因此有針對 unit 背後 repr 的特性進行分析的 work，例如 CITEUSPLEASE。 

\subsection{語音表徵的語音學分析}
在 HuBERT 出來之後，有一些研究像是 cite 等等，試圖探討對於語音表徵這樣語音模型的基礎進行各種從統計和語音學領域知識角度的分析，以期望能夠解釋為什麼模型可以擁有如此的表現。

此後，cite 等等作品則是從原先連續的表徵出發，開始往離散的量化向量，甚至是離散單元進行分析比對。雖然分析的切入角度可以相當多樣，例如 ABX、tsne 降維分群等等，但本次研究主要著重比對兩者之間在同一段語音序列上給予標籤的相關性，也就是以「偽標籤（pseudo-label）」的角度進行衡量。

\subsection{無文字（textless）語言模型}

這系列 textless 以 GSLM 為最主要代表作，旨在探討 unit 作為一種替代文字的方案。

本論文以 GSLM textless 採用的模型 units 為主要分析對象，企圖銜接兩者的脈絡，來佐證這些 unit 作為一種「類似或可替代文字的語音紀錄方式」在能夠發揮 LM 的特長背後，是否是基於符合語音學特徵帶來的，抑或有什麼其他特徵。

\section{衡量方式（metric）}

為了測量這些 unit 跟 phn 這類語言學 labels 之間的 correlation，我們需要先介紹本論文會探討的指標


\subsection{音位（phoneme）長條圖（bar chart）}
要先對 unit 做頻率上的統計
\subsection{純度（purity）}
phoneme purity: 每個cluster unit內的phoneme purity，代表此unit是否有phoneme代表性
cluster purity：每個phoneme對應到的cluster統計，若cluster purity低代表less liguistical meaning?（抄hubert paper）單一 phoneme 本身對應的 unit 的一致性。如果

\subsection{熵（entropy）和相互資訊（mutual information，MI）}
除了「最大」的對應關係，根據 Hubert 原先的 paper CITEME (hubert) 中的分析方式，我們也可以從 info theory 的角度，去探討「觀察到一個 unit 對於 label 不確定性的降低」來考慮 unit 本身提供了多少背後 phn 的資訊

\subsection{對齊（alignment）}
cluster是否保留segment資訊，不將不同phoneme合併
segmentation → 
怕說只是每個 frame 本身 unit 或 piece 可以跟 phoeneme 特徵相關，但放在連續的語音中被切得很碎/前後文相關的東西不知道有沒有抓到


\section{語音學分類（phone type）}

\subsection{簡介}

除了單一 phn 本身的特性以外，由於 phn 本身彼此不是完全獨立的，而是彼此之間就存在相似的特徵，可以分成幾個組別。因此，依照 CITEME (tanghao等三篇) 的分組方式，對英語的 phn 進行分類並合併比對數據，看看這些 unit 本身是否有 capture 到相似的發聲特徵，而不單純只是把 phn 分成約五十類完全獨立的標籤。（基於語音表徵本身就是 acoustic sisgnals 來的，應該 by nature 要可以對語音特徵分組吧？）

以下為各分組進行簡單介紹：

\subsubsection{consonants}

子音可以分成五類

\subsubsection{vowels}

母音在這邊為了簡單起見，會被分在一起？

\subsection{解釋意義}

\begin{itemize}
    \item 純度（purity）：換成type之後有何變化（關聯性更強？）
    \item 熵（entropy）（放直方圖解釋） --> phone type更明顯？
    \item 對齊（alignment）：是否減少segment資訊的保留（連續子音母音被合併？
\end{itemize}

\section{分析結果}

 
\subsection{基於各自音位的分析}
\subsection{基於語音學分類的分析}
\section{本章總結}

>>>>>>> b1b9a6a (Update on Overleaf.)
% \chapter{多個語音離散表徵組合與語音標記間的關係}



\section{動機}  % \section{簡介}
\section{相關研究}
\subsection{對語音離散表徵的分詞（tokenization）研究}
\section{分詞方法}
\section{衡量方式}
\subsection{字符（token）與音位之間的關係.}
\subsection{壓縮比率}
\section{分析結果}
\subsection{基於各自音位的分析}
\subsection{基於語音學分類的分析}
\section{應用在語音任務的實驗}
\subsection{語音辨識}
\subsubsection{實驗設定與資料集}
\subsubsection{實驗結果與其和分析數據間的關係}
\section{本章總結}

% \chapter{結論與展望}

\section{研究貢獻與討論}
\section{未來展望}

\newpage
\ % The empty page
\newpage

% !TeX root = ../thesis.tex

% !TeX root = ../../thesis.tex

\newcommand{\NOTEHERE}[1]{\colorbox{yellow}{\textcolor{blue}{
    \lparen\lparen\lparen}\textcolor{
        red}{#1}\textcolor{blue}{\rparen\rparen\rparen}}}

% newline
\newcommand{\PROBABLY}[1]{}
%% \newcommand{\PROBABLY}[1]{\colorbox{red}{\textcolor{green}{
%%     \lbrack\lbrack\lbrack}\textcolor{green}{#1}\textcolor{
%%         green}{\rbrack\rbrack\rbrack}}}

% \textcolor{green}{#1}
\newcommand{\FIXTHIS}[1]{\par\makebox[\textwidth]{
    \colorbox{pink}{\textcolor{green}{
        \lbrack\lbrack\lbrack}\textcolor{
            blue}{#1}\textcolor{green}{\rbrack\rbrack\rbrack}}}}

\newcommand{\mysetcounter}[2]{\setcounter{#1}{\the\numexpr#2 - 1\relax}}
\newcommand{\mychcnt}[1]{\mysetcounter{chapter}{#1}}

% % !TeX root = ../thesis.tex
\chapter{導論}
\section{研究動機}

　　語言是人與人彼此交流最主要的橋樑，而人們互相溝通最自然的方式便是透過說話的語音（speech）達成。人類往往是自幼就牙牙學語開始說話，直到已屆學齡左右才開始學習認字與書寫。雖然在這個資訊爆炸的時代，人們已經習慣以文字呈現的語言作為獲取資訊的主要媒介，但不論如何，各種書寫系統其背後承載的語言必定有語音的形式作為對應。更何況世界上現存大約七千多種  \cite{ethnologue}  語言中，絕大多數不見得存在成熟且普及的文字系統，卻無礙於這些語言被人們所熟悉和使用。因此，「語音」作為語言不可或缺的存在方式，了解它和研究它的價值自然不言而喻。

　　然而，相對於穩定、易於處理和保存的文字文本，語音訊號的變化萬千，蘊藏了大量從語者風格、表達內容到抑揚頓挫（韻律，prosody）等不同層次的訊息，使得對它的處理、研究相比之下複雜度與難度劇增。由於語音的這種特性，過往對於語言最有興趣的語言學家們，即便明白語音作為多數語言主體的事實，也不得不藉文字符號為依託進行探索。進入資訊化時代後，藉助電腦硬體等計算設備的幫助，從語料庫、計算語言學到自然語言處理等透過科技的力量發展語言處理技術的領域，頗長一段時間也是專注於文字的處理與分析。
而嘗試結合訊號處理發展的語音技術領域，當時則是透過語言學家對語言的領域知識，例如從 phoneme、morphology、語法等等用以刻劃人類語音和語言特性的概念，將之結合機器學習建立模型，開發技術以方便人們能以語音這種更靈活的媒介，更好的讓電腦、手機等科技工具可以更接近「直接溝通」的使用方式，便利人們的日常生活。

　　近年來，由於圖形處理器（graphics processing unit，GPU）等硬體平行運算技術的進步，深層學習（deep learning）快速崛起成為人工智慧的主流，有了此項機器學習的技術，模型的彈性能夠更好的萃取資料、更貼近的尋找資料背後的機制並進行預測，使得人們不再非得依賴大量費時費工的人類標注過程，進而使得利用大量語料庫發展語言技術，進一步推進語言科技發展成為可能。尤其在自監督學習（self-supervised learning）技術出現之後，深層學習模型可以依照人們給定的方向，更細緻的從大量未標注、相較容易取得的語音或文字的語料，找出其中的語音、語法及語義等等結構，形成帶有對人類語言有前所未見表現的基石模型（foundation model），是這個領域的一大里程碑。尤其在以處理文字為主體的自然語言處理領域，甚至出現了幾乎使人類真偽難辨的生成式模型，改變了人們生活的方方面面。

＿＿借鏡文字方面的成功經驗，語音處理領域的研究者們也開始嘗試將語言模型（language model）的概念套用於變化莫測的語音訊號之上，原先人們藉助訊號處理知識一直使用的各種語音訊號特徵（feature）也在自監督學習的架構之下，出現了許多模型從大量語音資料中得到的「語音表徵（speech representation）」，作為精煉語音資訊的另外一種新選擇開始廣泛被採用。然而，相比於文字符號的穩定與單純，語音的複雜性使得它處理起來會需要更大量的資料和運算資源來擷取其中不同層次的細節，而且作為物理訊號，語音還必須處理掉環境中的雜訊等干擾。為了從紛亂的聲音中提取出最重要的訊息，向量量化（vector quantization）的技巧因而經常被使用在語音 \cite{chorowski2019unsupervised, chen2023vector, zhao2023speech} 或影像的領域中。爾後， \cite{lakhotia2021generative}  基於模仿人類學習語言的過程，藉助諸如 CPC（\cite{oord2019representation} ）、HuBERT （\cite{hsu2021hubert} ）、wav2vec 2.0 （\cite{baevski2020wav2vec} ）等自監督學習模型的幫助，引入向量量化的技術，提出了「無文字（textless）」的學習架構，轉而以語音表徵量化後的「離散單元（discrete unit）」作為操作對象，企圖以單純大量的語音資料中訓練出一個不依賴文字的語言模型。此種學習架構的優勢在於在能保有利用大量未標注文字轉寫語音資料的同時，與連續表徵相比資訊的位元率（bit rate）利用更有效率、容易儲存、處理與傳輸，以及形式上更像文字的特性。自此之後有一系列  \NOTEHERE{補足 unit 相關的續作，例如英翻臺} 使用 discrete units 進行任務訓練的研究，一定程度的印證了unit 捕捉語音內容的有效性。

　　Unit 固然在編碼語音之上有不錯的效果，因此 empirically 被領域內的不少學者直接視為「機器自己學習出來的文字」，可是究竟這項 claim 本身有多符合實際狀況，有學者便嘗試重新從語言學的概念汲取靈感，對其進行語音學層面的分析析析析 \NOTEHERE{放唐顥等等的語音學研究}你知道吧？
　　\PROBABLY{你可能要併入連續的 speech repr}
　　\FIXTHIS{確認一下 acoustic piece 那邊的文線是不是也是如我那樣說的}
　　\FIXTHIS{確認研究動機要寫到哪裡}


\section{研究方向}

　　本論文旨在探討自監督學習模型的離散語音表徵與音位（phoneme）之間的關係。

\section{主要貢獻}
\section{章節安排}

本論文之章節安排如下：

% \chapter{導論}
\section{研究動機}

　　語言是人與人之間最主要的溝通方式，是我們對外界世界互動最主要的媒介。隨著時代的發展，在這個資訊爆炸的社會，人們接觸到語言、使用與互動的情形變得非常頻繁。然而，語言的多樣性在此卻成為一項阻礙不同人們交流的障礙，因此，發展語言科技為人們的互動提供協助，因而成為了一項不可避免的趨勢。藉助科技與知識的力量，人們可以透過不同語言的視角去獲取更多元的資訊，激發思考、開拓視野，因此諸如語音辨識、機器翻譯等等，一直都是人們熱衷於研究的議題。

    過往，為了達成此一遠大目標，機器學習工程師們與語言學家之間互相合作，期望藉由對領域知識的了解去建立模型，打造出創新的技術以應對語言的各種變化，來滿足人們溝通的需求。近年來，由於硬體平行運算技術的進步，深層學習（deep learning）快速崛起成為人工智慧的主流，有了此項機器學習的技術，模型的彈性能夠更好的萃取資料、更貼近的尋找資料背後的機制並進行預測，使得人們不再非得依賴大量費時費工的人類標註過程，進而使得利用大量語料庫發展語言技術，進一步推進語言科技發展成為可能。尤其在自監督學習（self-supervision）技術出現之後，深層學習模型可以依照人們給定的方向，更細緻的從大量未標注、相較容易取得的語音或文字的語料，找出其中的語音、語法及語義等等結構，形成帶有對人類語言有前所未見表現的基礎模型（foundation model），是這個領域的一大里程碑。尤其在以處理文字為主體的自然語言處理領域，甚至出現了幾乎使人類真偽難辨的生成式模型，改變了人們生活的方方面面。

    ……………………………………………………………………………… \\
    ………………………………………………………………………………………… \\
　　………………………………………………………………………………

    然而相較於穩定、易於處理的文字文本，語音訊號的變化複雜萬千，蘊藏了大量如內容、韻律、語者等等不同層次的訊息，使得處理的難度劇增。更何況，目前世界上大約七千多種語言中，絕大多數仍然沒有成熟且普及的文字系統。

    （先從語音表徵開始）

    因此，除了發展文字的模型以外，語音處理的技術更是必不可少的。
而在語音處理這邊，近期藉由自監督學習，提出了許多語音這邊的基礎模型。
% 這部分跟前面重疊到了
雖然

    因此在語音處理界，「無文字（textless）」的發展是相當吸引人的。

    （加上語言公平性問題？）

    近期由於 HuBERT 等等起來的 GSLM 等架構，已經在一定程度上做到完全不依賴文字轉寫、單純在大量的語音語料之上建立一定程度上媲美於「大型語言模型」的成果，甚至達成了閩南語和英語的互相對譯。此一里程碑大力的推動了語言科技的進展，有望推動藉助科技達成的降低語言障礙。

    ………………………………………………………………………………

    
    
   不過，即便語音的技術已經相當成熟，在追求模型表現的同時，語音與語言的技術開始與過往的人類對於語言學、語音的理解逐漸產生脫節。似乎在為了讓機器可以擁有良好表現的同時，理解模型如何運作似乎是被犧牲的必然。但人們在追求更好的模型表現的同時，有一群人開始注意到語音處理模型是否有可能抓取到人類語音中特有的、區別於一般音訊的特徵，並嘗試使用過往用來研究、歸類人類語音的方式，結合機器學習與統計學去解釋為什麼，並期求可以比較甚至改善機器模型在進行語音處理時的表現，不僅僅只是使用資料集本身的分數，而有更多更多元、更穩健的衡量標準。由於離散表徵在當今語音模型與語音處理技術已經愈來愈具備重要性，因此探討與分析為什麼語音離散表徵可以幫助下游任務的背後成因是相當重要的研究方向，其中一個驗證離散表徵能夠幫助模型處理語音訊號的方式，便是驗證其與音位（phoneme）之間的對應關係。

    我們想要知道的是，在離散單元推動語音處理發展的同時，它究竟與人類書寫和使用的文字還有多遠的差異，以及在使用上是否能夠達成如同文字的效果，仍舊是領域中尚待探討的議題。所以我們先看看，他是不是起碼符合語音學上，作為「捕捉語音內容」的基本特徵。

    （加上跟 LLM 的關係嗎？）

    也就是Phonology去討論這件事，看它跟聲學特徵跟語音到底有多像，從這個方向我們可以進而去改善說，也許我們可以去用不同的方式使用Unit，來進一步推進讓我們的Unwritten Language這些技術，能夠更好的跟，就是等於說進而發揮類似文字的效果，然後同時也能促進我們去知道，機器是怎麼樣去理解這些語音訊號，它可能看的是哪些部分。

    ……………………………………………………………………………… \\
    ………………………………………………………………………………………… \\
　　………………………………………………………………………………


\section{研究方向}

%%%     於是我們對這些離散單元進行了單一的分析
%%%     另外，因為人類的音素往往是超過一個音框的單位
%%%     因此推測離散單元單一本身會抓取到的是更加細緻的特徵
%%% 
%%%     那為了找出更加符合人類文字認知的表示符號
%%%     因此我們採用文字那邊常用的分詞器
%%%     // 可能要紀錄 rate 來證明這件事了
%%%     嘗試看看可否找出更接近人們認知信息率的符號，以其這種機器抓出
%%% 來的 fine structure 可以更往人類的文字靠近
%%% 
%%%     具體方法是，針對這些抓出來的 unit 跟 piece, 我們去算 purtiy,
%%% pmni, segmentation 這些指標，看看他們是不是真的跟人類 top down
%%% 標註出來的的 phoneme label 有統計上的一致性。
%%% 
%%%     最後我們嘗試做在 PR 或 ASR 上，驗證這個一致性不是僅僅在統計數
%%% 據上顯現，也是可以在應用上得到幫助的。

　　在近年，已經有不少相關的研究開始嘗試往將離散單元（discrete unit）作為除了連續表徵（continuous representation）以外，可以編碼（encode）語音訊號的另外一種方式。離散表徵（discrete representation）跟連續表徵相比，具有資訊更濃縮（位元率（bit rate）更低）因此更好儲存、處理與傳輸，以及形式上更像文字的特性。

    儘管離散表徵在語音社群（community）中常被當成一種類似文字的存在，另外有一些文獻則是將其當成連續表徵的精簡表示法。

    然而

    因此，我們藉由分析各種離散單元和人類理解語音最直接的處理層次：音位（phoneme）之間的關係，並將兩者進行各種統計上的序列比較，可以作為訓練大型語音基石模型（foundation model）的分詞（tokenization）基礎，選定最適合的表徵最為系統的輸入符記（token）。

    如 \cite{wu2023wav2seq} 等作品便是首先嘗試將離散單元進行 tokenization 後做進一步處理的，其後的


\section{主要貢獻}  % 有什麼

　　結果我們發現，藉由觀察這些 unit 並嘗試藉由分詞演算法找出更 high-level 的單位之後，我們觀察到這些機器學習 figure out 的「偽標記」一定程度上的符合了人類音素的特性，因此可以當成某種類似拼音文字的存在。當然這跟人類真正使用的拼音文字仍有距離，
% // 難道可以探討 text 是不是一種 speech tokenization 過程嗎？
因此雖然無法直接當成 exact 的文字使用，一些人類的標注還是需要的，不過透過 ML 我們已經可以盡量更有效的利用珍貴的標注資料，幫助那些尚不容易取得文字的語言發展語音語言科技，以協助保存他們的語言。
% // 問題是，有些語音任務並不需要抽內容嗎？


\section{章節安排}

　　由於本論文是以剖析既有的語音離散表徵為主軸，因此就相關研究方面需要從各角度入手，單獨成一章節。接著我們會從單一的離散單元，以及將單元視為像文字的字符（character）並進行分詞演算法兩種對語音離散單元處理的層次分別成章進行分析，最後將這些表徵嘗試做在語音的任務上，以驗證其具有一定的語音表徵能力，且能保留語音學的特徵。



%%%%%%%%%%%%%%%%%%%%%%%%%%%%%%%%%%%%%%%%%%%%%%%%%%%%%%%%%%%%%%%%%%%%%
%%%%%%%%%%%%%%%%%%%%%%%%%%%%%%%%%%%%%%%%%%%%%%%%%%%%%%%%%%%%%%%%%%%%%
%%%%%%%%%%%%%%%%%%%%%%%%%%%%%%%%%%%%%%%%%%%%%%%%%%%%%%%%%%%%%%%%%%%%%

\chapter{背景知識}

\section{深層類神經網路（deep neural network）}

深層類神經網路（Deep neural network） 是麥氏（McCulloch）在 1943 年提出 \cite{mcculloch1943logical} 仿生數學模型，旨在模擬生物神經系統的連結。

深層類神經網路是一個取法自生物神經連結的數學模型，其在計算認知神經科學中以連結派（connectionism）為主要代表，後在電腦科學與機器學習中有不同結構的進展。在此之後，基於其彈性與平行化的能力，能在 GPU 上面很有效率的進行運算並達到前所未有的效能，因此現在已經成為人工智慧發展的主流。

基於深層類神經網路的神經架構有 CNN、RNN、Transformer 等等，由於這些架構在語音與文字處理上都已經被廣泛使用，因此在下面分別介紹：


\subsection{卷積式（convolutional）類神經網路}

卷積式類神經網路一開始是在 \\cite{} 中提出，主要是鑑於影像中的局部性（locality），讓 NN 可以在

。

在語音中，因為

語音訊號的資訊是被呈現在時間維度上
，
因此
通常使用

一維

的
卷積式類神經網路
，
以捕捉
時間維度上的局部性特徵
，例如本研究特別探討的 phoneme、morpheme 等等。

% \subsection{遞迴式（recurrent）類神經網路}



\subsection{
    序列至序列（sequence-to-sequence）模型}

    


\subsection{專注（attention）機制}    

\subsection{轉換器（Transformer）}



\section{表徵（representation）學習}


\subsection{文字的語意表徵}

\subsection{語音特徵與表徵}

\section{語音基石模型與自監督式學習}

\subsection{自監督式學習}

\subsection{語音基石模型}

\subsection{離散單元}

\section{本章總結}

\setcounter{chapter}{1}
<<<<<<< HEAD
% !TeX root = ../thesis.tex
\mychcnt{3}\chapter{
單一語音離散表徵與語音標記的對應模式}
\section{動機}  % \section{簡介}
\section{相關研究}

在 HuBERT 出來之後，

\subsection{語音表徵的語音學分析}
\section{衡量方式}
\subsection{音位（phoneme）長條圖（bar chart）}
\subsection{純度（purity）、熵（entropy）}
\subsection{對齊（alignment）}
\section{語音學分類}
\section{分析結果}
\subsection{基於各自音位的分析}
\subsection{基於語音學分類的分析}
\section{本章總結}

% ============================================================ %
% ============================================================ %

\mychcnt{2}\chapter{
背景知識}
\section{深層類神經網路（deep neural network）}

　　深層類神經網路（Deep neural network） 是麥氏（McCulloch）在 1943 年提出 \cite{mcculloch1943logical}，取法自生物神經連結的計算模型，旨在模擬生物神經系統的連結，以模仿生物的各項功能，進而透過機器學習的最佳化演算法，使得整個模型能夠藉由資料去貼合理想的函數，以達成應用或工程上所需要的各種任務。 。

    以發展此模型為主軸的心理學流派，在計算認知神經科學中被稱為「連結派（connectionism）」，其後因為該網路的彈性與平行化的能力，和諸如圖形處理器（graph processing unit，GPU）等硬體裝置能夠 最好的利用。 
    并能夠更 好的描述資料分佈、達到前所未有的效能 ，因此近年在電腦科學的機器學習領域中獲得重大進展，並因此現已成為人工智慧發展的主流。

    基於深層類神經網路的神經架構有 CNN、RNN、Transformer 等等，由於這些架構在語音與文字處理上都已經被廣泛使用，因此在下面分別介紹：

\subsection{卷積式（convolutional）類神經網路}

　　卷積式類神經網路一開始是在 cite 中提出，主要是鑑於影像中的局部性（locality），讓 NN 可以在。在語音中，因為語音訊號的資訊是被呈現在時間維度上，因此通常使用一維的卷積式類神經網路，以捕捉時間維度上的局部性特徵，例如本研究特別探討的 phoneme、morpheme 等等。

\subsection{遞迴式（recurrent）類神經網路}

\subsection{
    序列至序列（sequence-to-sequence）模型}
    
    由于許多實際上的資料都是2個序 之間互相配對的關系 此類的資料包含語音文字。 信號等等，都是以時間軸為主要 演變方向的資料。 因此 有一類模型。 會被以序列制序列的模式進行訓練。 旨在模擬輸入與輸出序列之間的變化與相依關係（dependency）。

此類模型一般的架構是由一個編碼器和解碼器構成 其中編碼器是將輸入訊號借由內部表征進行編碼。 依據每個時間點輸入訊號的順序來改變其內部表征的狀態 接著將最後一個時間點的表徵作為整個序列的特征傳遞給解碼器進行輸出訊號 生成。 
    
\subsection{專注（attention）機制}
    
    原本的序列自序列模型本身。 需要讓解碼器單純透過最後一個時間 點。 的表征資訊來完全儲存輸入序列的一切資訊 以工解碼器判斷。 并生成輸出序列。 然而，由于 單就最後一個向量進行判斷對於解碼器而言過於不易。因此。 盧氏提出。 在編碼器中對輸入序列的不同時間點進行注意力機制 亦即讓解碼器可以根據當下 所需要輸出的 內容判斷應該要重新對輸入序列的哪些部份 進行更多的加權 。

    

\subsection{轉換器（Transformer）}

其後，由瓦氏（Vaswani） cite 提出的 論文中 提出了一個完全由注意機制。 所構成的序列自序列模型。 原先該模型適用於解決機器翻譯。 的問題。
由於其能夠高度平行化的特性， 日後在 自然源處理和語音處理，甚至到電腦視覺領域等近乎整個深層學習的領域都被廣泛的應用。  


\section{表徵（representation）學習}

\subsection{文字的語意表徵}

word2vec 是 mikolov \CITEME 最早提出跟對於文字進行語義表徵的 work。在其後

\subsection{語音特徵與表徵}


\section{語音基石模型與自監督式學習}

\subsection{自監督式學習}

\subsection{語音基石模型}

\subsection{離散單元}


\section{本章總結}



% % !TeX root = ../thesis.tex
\mychcnt{3}\chapter{
    單一語音離散表徵與語音標記的對應模式}
\section{動機}  % \section{簡介}
\section{相關研究}

在 HuBERT 出來之後，

\subsection{語音表徵的語音學分析}
\section{衡量方式}
\subsection{音位（phoneme）長條圖（bar chart）}
\subsection{純度（purity）、熵（entropy）}
\subsection{對齊（alignment）}
\section{語音學分類}
\section{分析結果}
\subsection{基於各自音位的分析}
\subsection{基於語音學分類的分析}
\section{本章總結}

\mychcnt{2}\chapter{背景知識}
\section{真全部深層類神經網路（deep neural network）}

　　深層類神經網路（Deep neural network） 是麥氏（McCulloch）在 1943 年提出 \cite{mcculloch1943logical}，取法自生物神經連結的計算模型，旨在模擬生物神經系統的連結，以模仿生物的各項功能，進而透過機器學習的最佳化演算法，使得整個模型能夠藉由資料去貼合理想的函數，以達成應用或工程上所需要的各種任務。 。

    以發展此模型為主軸的心理學流派，在計算認知神經科學中被稱為「連結派（connectionism）」，其後因為該網路的彈性與平行化的能力，和諸如圖形處理器（graph processing unit，GPU）等硬體裝置能夠 最好的利用。 
    并能夠更 好的描述資料分佈、達到前所未有的效能 ，因此近年在電腦科學的機器學習領域中獲得重大進展，並因此現已成為人工智慧發展的主流。

    基於深層類神經網路的神經架構有 CNN、RNN、Transformer 等等，由於這些架構在語音與文字處理上都已經被廣泛使用，因此在下面分別介紹：

\subsection{卷積式（convolutional）類神經網路}

　　卷積式類神經網路一開始是在 cite 中提出，主要是鑑於影像中的局部性（locality），讓 NN 可以在。在語音中，因為語音訊號的資訊是被呈現在時間維度上，因此通常使用一維的卷積式類神經網路，以捕捉時間維度上的局部性特徵，例如本研究特別探討的 phoneme、morpheme 等等。

\subsection{遞迴式（recurrent）類神經網路}

\subsection{
    序列至序列（sequence-to-sequence）模型}
    
    由于許多實際上的資料都是2個序 之間互相配對的關系 此類的資料包含語音文字。 信號等等，都是以時間軸為主要 演變方向的資料。 因此 有一類模型。 會被以序列制序列的模式進行訓練。 旨在模擬輸入與輸出序列之間的變化與相依關係（dependency）。

此類模型一般的架構是由一個編碼器和解碼器構成 其中編碼器是將輸入訊號借由內部表征進行編碼。 依據每個時間點輸入訊號的順序來改變其內部表征的狀態 接著將最後一個時間點的表徵作為整個序列的特征傳遞給解碼器進行輸出訊號 生成。 
    
\subsection{專注（attention）機制}
    
    原本的序列自序列模型本身。 需要讓解碼器單純透過最後一個時間 點。 的表征資訊來完全儲存輸入序列的一切資訊 以工解碼器判斷。 并生成輸出序列。 然而，由于 單就最後一個向量進行判斷對於解碼器而言過於不易。因此。 盧氏提出。 在編碼器中對輸入序列的不同時間點進行注意力機制 亦即讓解碼器可以根據當下 所需要輸出的 內容判斷應該要重新對輸入序列的哪些部份 進行更多的加權 。

    

\subsection{轉換器（Transformer）}

其後，由瓦氏（Vaswani） cite 提出的 論文中 提出了一個完全由注意機制。 所構成的序列自序列模型。 原先該模型適用於解決機器翻譯。 的問題。
由於其能夠高度平行化的特性， 日後在 自然源處理和語音處理，甚至到電腦視覺領域等近乎整個深層學習的領域都被廣泛的應用。  


\section{表徵（representation）學習}

\subsection{文字的語意表徵}



\subsection{語音特徵與表徵}


\section{語音基石模型與自監督式學習}

\subsection{自監督式學習}

\subsection{語音基石模型}

\subsection{離散單元}


\section{本章總結}



% % !TeX root = ../thesis.tex
\mychcnt{3}\chapter{
單一語音離散表徵與語音標記的對應模式}
\section{動機}  % \section{簡介}
\section{相關研究}

在 HuBERT 出來之後，

\subsection{語音表徵的語音學分析}
\section{衡量方式}
\subsection{音位（phoneme）長條圖（bar chart）}
\subsection{純度（purity）、熵（entropy）}
\subsection{對齊（alignment）}
\section{語音學分類}
\section{分析結果}
\subsection{基於各自音位的分析}
\subsection{基於語音學分類的分析}
\section{本章總結}

% ============================================================ %
% ============================================================ %

\mychcnt{2}\chapter{
背景知識}
\section{深層類神經網路（deep neural network）}

　　深層類神經網路（Deep neural network） 是麥氏（McCulloch）在 1943 年提出 \cite{mcculloch1943logical}，取法自生物神經連結的計算模型，旨在模擬生物神經系統的連結，以模仿生物的各項功能，進而透過機器學習的最佳化演算法，使得整個模型能夠藉由資料去貼合理想的函數，以達成應用或工程上所需要的各種任務。 。

    以發展此模型為主軸的心理學流派，在計算認知神經科學中被稱為「連結派（connectionism）」，其後因為該網路的彈性與平行化的能力，和諸如圖形處理器（graph processing unit，GPU）等硬體裝置能夠 最好的利用。 
    并能夠更 好的描述資料分佈、達到前所未有的效能 ，因此近年在電腦科學的機器學習領域中獲得重大進展，並因此現已成為人工智慧發展的主流。

    基於深層類神經網路的神經架構有 CNN、RNN、Transformer 等等，由於這些架構在語音與文字處理上都已經被廣泛使用，因此在下面分別介紹：

\subsection{卷積式（convolutional）類神經網路}

　　卷積式類神經網路一開始是在 cite 中提出，主要是鑑於影像中的局部性（locality），讓 NN 可以在。在語音中，因為語音訊號的資訊是被呈現在時間維度上，因此通常使用一維的卷積式類神經網路，以捕捉時間維度上的局部性特徵，例如本研究特別探討的 phoneme、morpheme 等等。

\subsection{遞迴式（recurrent）類神經網路}

\subsection{
    序列至序列（sequence-to-sequence）模型}
    
    由于許多實際上的資料都是2個序 之間互相配對的關系 此類的資料包含語音文字。 信號等等，都是以時間軸為主要 演變方向的資料。 因此 有一類模型。 會被以序列制序列的模式進行訓練。 旨在模擬輸入與輸出序列之間的變化與相依關係（dependency）。

此類模型一般的架構是由一個編碼器和解碼器構成 其中編碼器是將輸入訊號借由內部表征進行編碼。 依據每個時間點輸入訊號的順序來改變其內部表征的狀態 接著將最後一個時間點的表徵作為整個序列的特征傳遞給解碼器進行輸出訊號 生成。 
    
\subsection{專注（attention）機制}
    
    原本的序列自序列模型本身。 需要讓解碼器單純透過最後一個時間 點。 的表征資訊來完全儲存輸入序列的一切資訊 以工解碼器判斷。 并生成輸出序列。 然而，由于 單就最後一個向量進行判斷對於解碼器而言過於不易。因此。 盧氏提出。 在編碼器中對輸入序列的不同時間點進行注意力機制 亦即讓解碼器可以根據當下 所需要輸出的 內容判斷應該要重新對輸入序列的哪些部份 進行更多的加權 。

    

\subsection{轉換器（Transformer）}

其後，由瓦氏（Vaswani） cite 提出的 論文中 提出了一個完全由注意機制。 所構成的序列自序列模型。 原先該模型適用於解決機器翻譯。 的問題。
由於其能夠高度平行化的特性， 日後在 自然源處理和語音處理，甚至到電腦視覺領域等近乎整個深層學習的領域都被廣泛的應用。  


\section{表徵（representation）學習}

\subsection{文字的語意表徵}

word2vec 是 mikolov \CITEME 最早提出跟對於文字進行語義表徵的 work。在其後

\subsection{語音特徵與表徵}


\section{語音基石模型與自監督式學習}

\subsection{自監督式學習}

\subsection{語音基石模型}

\subsection{離散單元}


\section{本章總結}



% % !TeX root = ../../thesis.tex

% % !TeX root = ../thesis.tex

\chapter{
單一語音離散表徵與語音標記的對應模式}
\section{動機簡介}  % \section{簡介}
由於 HuBERT 之後，unit 的使用很廣泛，因此為了研究 unit 本身為什麼會被如此適當的可以讓模型視為文字對語音資料進行訓練，我們先從離散表徵本身的特徵分析起。 

% // 從為了驗證它是不是文字，改成分析內部有沒有精細結構
% ?? 要不要去分析有沒有 dseduplicate 這件事對於影響精細結構的關係，以及 rate 等等的差異呢？

\section{相關研究}
近期已經有多項相關的研究，嘗試在 SSL 這麼厲害的表現之後找原因，因此有針對 unit 背後 repr 的特性進行分析的 work，例如 CITEUSPLEASE。 

\subsection{語音表徵的語音學分析}
在 HuBERT 出來之後，有一些研究像是 cite 等等，試圖探討對於語音表徵這樣語音模型的基礎進行各種從統計和語音學領域知識角度的分析，以期望能夠解釋為什麼模型可以擁有如此的表現。

此後，cite 等等作品則是從原先連續的表徵出發，開始往離散的量化向量，甚至是離散單元進行分析比對。雖然分析的切入角度可以相當多樣，例如 ABX、tsne 降維分群等等，但本次研究主要著重比對兩者之間在同一段語音序列上給予標籤的相關性，也就是以「偽標籤（pseudo-label）」的角度進行衡量。

\subsection{無文字（textless）語言模型}

這系列 textless 以 GSLM 為最主要代表作，旨在探討 unit 作為一種替代文字的方案。

本論文以 GSLM textless 採用的模型 units 為主要分析對象，企圖銜接兩者的脈絡，來佐證這些 unit 作為一種「類似或可替代文字的語音紀錄方式」在能夠發揮 LM 的特長背後，是否是基於符合語音學特徵帶來的，抑或有什麼其他特徵。

\section{衡量方式（metric）}

為了測量這些 unit 跟 phn 這類語言學 labels 之間的 correlation，我們需要先介紹本論文會探討的指標


\subsection{音位（phoneme）長條圖（bar chart）}
要先對 unit 做頻率上的統計
\subsection{純度（purity）}
phoneme purity: 每個cluster unit內的phoneme purity，代表此unit是否有phoneme代表性
cluster purity：每個phoneme對應到的cluster統計，若cluster purity低代表less liguistical meaning?（抄hubert paper）單一 phoneme 本身對應的 unit 的一致性。如果

\subsection{熵（entropy）和相互資訊（mutual information，MI）}
除了「最大」的對應關係，根據 Hubert 原先的 paper CITEME (hubert) 中的分析方式，我們也可以從 info theory 的角度，去探討「觀察到一個 unit 對於 label 不確定性的降低」來考慮 unit 本身提供了多少背後 phn 的資訊

\subsection{對齊（alignment）}
cluster是否保留segment資訊，不將不同phoneme合併
segmentation → 
怕說只是每個 frame 本身 unit 或 piece 可以跟 phoeneme 特徵相關，但放在連續的語音中被切得很碎/前後文相關的東西不知道有沒有抓到


\section{語音學分類（phone type）}

\subsection{簡介}

除了單一 phn 本身的特性以外，由於 phn 本身彼此不是完全獨立的，而是彼此之間就存在相似的特徵，可以分成幾個組別。因此，依照 CITEME (tanghao等三篇) 的分組方式，對英語的 phn 進行分類並合併比對數據，看看這些 unit 本身是否有 capture 到相似的發聲特徵，而不單純只是把 phn 分成約五十類完全獨立的標籤。（基於語音表徵本身就是 acoustic sisgnals 來的，應該 by nature 要可以對語音特徵分組吧？）

以下為各分組進行簡單介紹：

\subsubsection{consonants}

子音可以分成五類

\subsubsection{vowels}

母音在這邊為了簡單起見，會被分在一起？

\subsection{解釋意義}

\begin{itemize}
    \item 純度（purity）：換成type之後有何變化（關聯性更強？）
    \item 熵（entropy）（放直方圖解釋） --> phone type更明顯？
    \item 對齊（alignment）：是否減少segment資訊的保留（連續子音母音被合併？
\end{itemize}

\section{分析結果}

 
\subsection{基於各自音位的分析}
\subsection{基於語音學分類的分析}
\section{本章總結}

=======
% !TeX root = ../../thesis.tex

% !TeX root = ../thesis.tex

\chapter{
單一語音離散表徵與語音標記的對應模式}
\section{動機簡介}  % \section{簡介}
由於 HuBERT 之後，unit 的使用很廣泛，因此為了研究 unit 本身為什麼會被如此適當的可以讓模型視為文字對語音資料進行訓練，我們先從離散表徵本身的特徵分析起。 

% // 從為了驗證它是不是文字，改成分析內部有沒有精細結構
% ?? 要不要去分析有沒有 dseduplicate 這件事對於影響精細結構的關係，以及 rate 等等的差異呢？

\section{相關研究}
近期已經有多項相關的研究，嘗試在 SSL 這麼厲害的表現之後找原因，因此有針對 unit 背後 repr 的特性進行分析的 work，例如 CITEUSPLEASE。 

\subsection{語音表徵的語音學分析}
在 HuBERT 出來之後，有一些研究像是 cite 等等，試圖探討對於語音表徵這樣語音模型的基礎進行各種從統計和語音學領域知識角度的分析，以期望能夠解釋為什麼模型可以擁有如此的表現。

此後，cite 等等作品則是從原先連續的表徵出發，開始往離散的量化向量，甚至是離散單元進行分析比對。雖然分析的切入角度可以相當多樣，例如 ABX、tsne 降維分群等等，但本次研究主要著重比對兩者之間在同一段語音序列上給予標籤的相關性，也就是以「偽標籤（pseudo-label）」的角度進行衡量。

\subsection{無文字（textless）語言模型}

這系列 textless 以 GSLM 為最主要代表作，旨在探討 unit 作為一種替代文字的方案。

本論文以 GSLM textless 採用的模型 units 為主要分析對象，企圖銜接兩者的脈絡，來佐證這些 unit 作為一種「類似或可替代文字的語音紀錄方式」在能夠發揮 LM 的特長背後，是否是基於符合語音學特徵帶來的，抑或有什麼其他特徵。

\section{衡量方式（metric）}

為了測量這些 unit 跟 phn 這類語言學 labels 之間的 correlation，我們需要先介紹本論文會探討的指標


\subsection{音位（phoneme）長條圖（bar chart）}
要先對 unit 做頻率上的統計
\subsection{純度（purity）}
phoneme purity: 每個cluster unit內的phoneme purity，代表此unit是否有phoneme代表性
cluster purity：每個phoneme對應到的cluster統計，若cluster purity低代表less liguistical meaning?（抄hubert paper）單一 phoneme 本身對應的 unit 的一致性。如果

\subsection{熵（entropy）和相互資訊（mutual information，MI）}
除了「最大」的對應關係，根據 Hubert 原先的 paper CITEME (hubert) 中的分析方式，我們也可以從 info theory 的角度，去探討「觀察到一個 unit 對於 label 不確定性的降低」來考慮 unit 本身提供了多少背後 phn 的資訊

\subsection{對齊（alignment）}
cluster是否保留segment資訊，不將不同phoneme合併
segmentation → 
怕說只是每個 frame 本身 unit 或 piece 可以跟 phoeneme 特徵相關，但放在連續的語音中被切得很碎/前後文相關的東西不知道有沒有抓到


\section{語音學分類（phone type）}

\subsection{簡介}

除了單一 phn 本身的特性以外，由於 phn 本身彼此不是完全獨立的，而是彼此之間就存在相似的特徵，可以分成幾個組別。因此，依照 CITEME (tanghao等三篇) 的分組方式，對英語的 phn 進行分類並合併比對數據，看看這些 unit 本身是否有 capture 到相似的發聲特徵，而不單純只是把 phn 分成約五十類完全獨立的標籤。（基於語音表徵本身就是 acoustic sisgnals 來的，應該 by nature 要可以對語音特徵分組吧？）

以下為各分組進行簡單介紹：

\subsubsection{consonants}

子音可以分成五類

\subsubsection{vowels}

母音在這邊為了簡單起見，會被分在一起？

\subsection{解釋意義}

\begin{itemize}
    \item 純度（purity）：換成type之後有何變化（關聯性更強？）
    \item 熵（entropy）（放直方圖解釋） --> phone type更明顯？
    \item 對齊（alignment）：是否減少segment資訊的保留（連續子音母音被合併？
\end{itemize}

\section{分析結果}

 
\subsection{基於各自音位的分析}
\subsection{基於語音學分類的分析}
\section{本章總結}

>>>>>>> b1b9a6a (Update on Overleaf.)
% \chapter{多個語音離散表徵組合與語音標記間的關係}



\section{動機}  % \section{簡介}
\section{相關研究}
\subsection{對語音離散表徵的分詞（tokenization）研究}
\section{分詞方法}
\section{衡量方式}
\subsection{字符（token）與音位之間的關係.}
\subsection{壓縮比率}
\section{分析結果}
\subsection{基於各自音位的分析}
\subsection{基於語音學分類的分析}
\section{應用在語音任務的實驗}
\subsection{語音辨識}
\subsubsection{實驗設定與資料集}
\subsubsection{實驗結果與其和分析數據間的關係}
\section{本章總結}

% \chapter{結論與展望}

\section{研究貢獻與討論}
\section{未來展望}

\newpage
\ % The empty page
\newpage

% % this file is encoded in utf-8
% v3.0 (Jun. 11, 2019)

%%% 參考文獻
\newpage
\phantomsection % for hyperref to register this
\addcontentsline{toc}{chapter}{\nameRef}
\renewcommand{\bibname}{\protect\makebox[5cm][s]{\nameRef}}
%  \makebox{} is fragile; need protect
\bibliographystyle{IEEEtran}  % 使用 IEEE Trans 期刊格式
\bibliography{thesis}


%%% 附錄
%\input{my_appendix.tex}

%%% 自傳
%\newpage
%\chapter*{\protect\makebox[5cm][s]{\nameVita}} % \makebox{} is fragile; need protect
%\phantomsection % for hyperref to register this
%\addcontentsline{toc}{chapter}{\nameVita}
%\input{my_vita.tex}



% <<<
\newpage
\phantomsection % for hyperref to register this
\addcontentsline{toc}{chapter}{\nameRef}
\renewcommand{\bibname}{\protect\makebox[3.5cm][s]{\nameRef}}  % fixed by JEFF
% \renewcommand{\bibname}{\protect\makebox[5cm][s]{\nameRef}}  % original
%  \makebox{} is fragile; need protect
\bibliographystyle{IEEEtran}  % 使用 IEEE Trans 期刊格式
\bibliography{thesis}

% >>>

\end{document}
